% VLC -- PPM-Signalgenerator (Analog Discovery 2)
% Dieses Dokument kann mit % VLC -- PPM-Signalgenerator (Analog Discovery 2)
% Dieses Dokument kann mit % VLC -- PPM-Signalgenerator (Analog Discovery 2)
% Dieses Dokument kann mit % VLC -- PPM-Signalgenerator (Analog Discovery 2)
% Dieses Dokument kann mit \input{README} in ein LaTeX-Hauptdokument eingebunden werden.
%
% Benoetigt folgende Pakete im Hauptdokument:
%   \usepackage[T1]{fontenc}
%   \usepackage[utf8]{inputenc}
%   \usepackage{booktabs}
%   \usepackage{listings}
%   \usepackage{hyperref}
%   \usepackage{enumitem}
%   \usepackage{microtype}

\section*{VLC -- PPM-Signalgenerator (Analog Discovery 2)}

Ein Python-basierter Pulspositionsmodulations-Signalgenerator (PPM),
der den \textbf{Digilent Analog Discovery 2} als Hardware-Ausgabegeraet verwendet.

% ------------------------------------------------------------------
\subsection*{Uebersicht}

Dieses Programm erzeugt ein PPM-Signal (Pulspositionsmodulation) aus einer
vom Benutzer eingegebenen binaeren Bitfolge und gibt es in Echtzeit ueber den
\textbf{Arbitrary Waveform Generator (AWG)} des Analog Discovery 2 aus.

Der Einstiegspunkt des Programms ist \textbf{\texttt{main\_PPM.py}} im
Verzeichnis \texttt{VLC\_Waveform SDK/}.

% ------------------------------------------------------------------
\subsection*{Funktionsweise von PPM}

Bei PPM wird jedes Bit durch einen Impuls dargestellt, der an einer bestimmten
Position innerhalb eines festen Zeitfensters platziert wird:

\begin{center}
\begin{tabular}{ll}
\toprule
\textbf{Bit} & \textbf{Impulsposition} \\
\midrule
\texttt{0} & Impuls am \textbf{Anfang} des Zeitfensters \\
\texttt{1} & Impuls in der \textbf{Mitte} des Zeitfensters \\
\bottomrule
\end{tabular}
\end{center}

Die Standardspannung betraegt \texttt{+1 V} fuer Impulse und \texttt{0 V}
fuer den Ruhepegel. Zeitverhalten und Amplitude sind vollstaendig in
\texttt{config.py} konfigurierbar.

% ------------------------------------------------------------------
\subsection*{Hardwareanforderungen}

\begin{itemize}
  \item \textbf{Digilent Analog Discovery 2} (AD2) -- USB-Multifunktionsinstrument
  \begin{itemize}
    \item Wird als Arbitrary Waveform Generator (AWG) verwendet, um das
          PPM-Signal am \textbf{Analogausgang Kanal 1 (W1)} auszugeben
  \end{itemize}
  \item Ein USB-Kabel zur Verbindung des AD2 mit dem PC
\end{itemize}

% ------------------------------------------------------------------
\subsection*{Softwarevoraussetzungen}

\subsubsection*{1. WaveForms (Digilent SDK)}

Das Programm benoetigt die \textbf{DWF-Bibliothek (Device WaveForms)}
(\texttt{dwf.dll} unter Windows, \texttt{libdwf.so} unter Linux,
\texttt{libdwf.dylib} unter macOS).

WaveForms kann von der Digilent-Website heruntergeladen und installiert werden:\\
\url{https://digilent.com/reference/software/waveforms/waveforms-3/start}

\textbf{Wichtig:} Die WaveForms-Desktopanwendung muss geschlossen sein, bevor
ein Python-Skript ausgefuehrt wird, da WaveForms und ein Python-Skript das
Geraet nicht gleichzeitig steuern koennen.

\subsubsection*{2. Python 3}

Python 3 installieren (Version 3.7 oder hoeher empfohlen).

\subsubsection*{3. Python-Pakete}

Das benoetigt Paket installieren:

\begin{lstlisting}[language=bash]
pip install matplotlib
\end{lstlisting}

\texttt{ctypes} und \texttt{time} sind Bestandteil der Python-Standardbibliothek
und muessen nicht separat installiert werden.

% ------------------------------------------------------------------
\subsection*{Dateistruktur}

\begin{lstlisting}
VLC_Waveform SDK/
+-- main_PPM.py       <- Hauptskript -- diese Datei ausfuehren
+-- PPM.py            <- Logik zur PPM-Signalerzeugung
+-- config.py         <- Alle konfigurierbaren Parameter und Geraeteeinrichtung
+-- wavegen.py        <- Low-Level-Hilfsfunktionen fuer den Waveformgenerator
+-- device.py         <- Hilfsfunktionen fuer Geraeteverbindung und Fehlerbehandlung
+-- dwfconstants.py   <- DWF-Bibliothekskonstanten
+-- main.py           <- Eigenstaendige Demo: Rechteckwelle (separater Test)
+-- waveform_data.csv <- (optional) Exportierte Wellenformdaten
\end{lstlisting}

% ------------------------------------------------------------------
\subsection*{Konfiguration}

Vor dem Starten des Programms sollten die Parameter in \textbf{\texttt{config.py}}
geprueft und bei Bedarf angepasst werden:

\begin{center}
\begin{tabular}{lll}
\toprule
\textbf{Parameter} & \textbf{Standardwert} & \textbf{Beschreibung} \\
\midrule
\texttt{DESIRED\_PEAK\_VOLTAGE} & \texttt{1.0} V & Impulsamplitude \\
\texttt{REST\_LEVEL} & \texttt{0.0} V & Spannung waehrend des Ruhepegels \\
\texttt{BIT\_DURATION} & \texttt{1e-6} s (1\,\textmu s) & Dauer eines Bit-Zeitfensters \\
\texttt{PULSE\_WIDTH} & \texttt{BIT\_DURATION / 2} & Impulsbreite innerhalb eines Zeitfensters \\
\texttt{SAMPLES\_PER\_BIT\_GEN} & \texttt{2048} & AWG-Abtastwerte pro Bit \\
\texttt{LOOP\_CONTINUOUSLY} & \texttt{True} & Wellenform wiederholen (\texttt{True}) oder einmalig (\texttt{False}) \\
\texttt{CHANNEL\_AWG} & \texttt{0} & AWG-Ausgangskanal (0 = W1) \\
\texttt{save\_to\_CSV} & \texttt{False} & Abtastwerte in \texttt{waveform\_data.csv} speichern \\
\texttt{diagnostic} & \texttt{False} & Signaldetails in der Konsole ausgeben \\
\bottomrule
\end{tabular}
\end{center}

% ------------------------------------------------------------------
\subsection*{Programm starten}

\begin{enumerate}
  \item Den Analog Discovery 2 per USB mit dem Computer verbinden.
  \item Die WaveForms-Desktopanwendung schliessen, falls sie geoeffnet ist.
  \item Ein Terminal oeffnen und in den SDK-Ordner wechseln:
\begin{lstlisting}[language=bash]
cd "VLC_Waveform SDK"
\end{lstlisting}
  \item Das Hauptskript ausfuehren:
\begin{lstlisting}[language=bash]
python main_PPM.py
\end{lstlisting}
  Bei Erfolg erscheint folgende Ausgabe:
\begin{lstlisting}
Device opened successfully.
Bitfolge in Zahlen eintragen...
\end{lstlisting}
\end{enumerate}

% ------------------------------------------------------------------
\subsection*{Programm verwenden}

Nach dem Start wechselt das Programm in eine interaktive Schleife:

\begin{enumerate}
  \item Eine binaere Bitfolge eingeben, wenn der Prompt erscheint --
        eine Folge aus \texttt{0} und \texttt{1}, zum Beispiel:
\begin{lstlisting}
Bitfolge in Zahlen eintragen...01101001
\end{lstlisting}
  \item Das Programm fuehrt anschliessend folgende Schritte aus:
  \begin{itemize}
    \item Es erzeugt die entsprechende PPM-Wellenform
    \item Es gibt das Signal ueber den AWG-Kanal (W1) des Analog Discovery 2 aus
    \item Es wartet, bis die Signaldauer abgelaufen ist, und fragt dann erneut
          nach einer Eingabe
  \end{itemize}
  \item Den Vorgang mit einer neuen Bitfolge beliebig oft wiederholen.
  \item Das Programm mit \texttt{Ctrl+C} beenden. Der AWG-Ausgang wird gestoppt
        und das Geraet ordnungsgemaess getrennt.
\end{enumerate}

\textbf{Hinweis:} Das erzeugte Signal kann mit dem Oszilloskopeingang des
Analog Discovery 2 (Kanal 1 / 1+) und dem WaveForms-Scope-Instrument beobachtet
werden -- jedoch erst nach dem Beenden von \texttt{main\_PPM.py}, da jeweils nur
eine Anwendung das Geraet steuern kann.

% ------------------------------------------------------------------
\subsection*{Fehlerbehebung}

\begin{center}
\begin{tabular}{p{0.38\textwidth}p{0.55\textwidth}}
\toprule
\textbf{Problem} & \textbf{Loesung} \\
\midrule
\texttt{Failed to open device} &
  Sicherstellen, dass der AD2 per USB verbunden und WaveForms geschlossen ist \\
\addlinespace
\texttt{Could not load dwf library} &
  WaveForms von Digilent installieren; sicherstellen, dass die DWF-Bibliothek
  im System-PATH vorhanden ist \\
\addlinespace
\texttt{No device detected} &
  Einen anderen USB-Anschluss oder ein anderes Kabel verwenden;
  Geraetemanager bzw.\ \texttt{lsusb} pruefen \\
\addlinespace
Skript haengt oder keine Ausgabe &
  \texttt{BIT\_DURATION} und Bitfolgenlange pruefen; diese sollten nicht
  zu gross sein \\
\addlinespace
\texttt{ModuleNotFoundError: matplotlib} &
  \texttt{pip install matplotlib} ausfuehren \\
\bottomrule
\end{tabular}
\end{center}
 in ein LaTeX-Hauptdokument eingebunden werden.
%
% Benoetigt folgende Pakete im Hauptdokument:
%   \usepackage[T1]{fontenc}
%   \usepackage[utf8]{inputenc}
%   \usepackage{booktabs}
%   \usepackage{listings}
%   \usepackage{hyperref}
%   \usepackage{enumitem}
%   \usepackage{microtype}

\section*{VLC -- PPM-Signalgenerator (Analog Discovery 2)}

Ein Python-basierter Pulspositionsmodulations-Signalgenerator (PPM),
der den \textbf{Digilent Analog Discovery 2} als Hardware-Ausgabegeraet verwendet.

% ------------------------------------------------------------------
\subsection*{Uebersicht}

Dieses Programm erzeugt ein PPM-Signal (Pulspositionsmodulation) aus einer
vom Benutzer eingegebenen binaeren Bitfolge und gibt es in Echtzeit ueber den
\textbf{Arbitrary Waveform Generator (AWG)} des Analog Discovery 2 aus.

Der Einstiegspunkt des Programms ist \textbf{\texttt{main\_PPM.py}} im
Verzeichnis \texttt{VLC\_Waveform SDK/}.

% ------------------------------------------------------------------
\subsection*{Funktionsweise von PPM}

Bei PPM wird jedes Bit durch einen Impuls dargestellt, der an einer bestimmten
Position innerhalb eines festen Zeitfensters platziert wird:

\begin{center}
\begin{tabular}{ll}
\toprule
\textbf{Bit} & \textbf{Impulsposition} \\
\midrule
\texttt{0} & Impuls am \textbf{Anfang} des Zeitfensters \\
\texttt{1} & Impuls in der \textbf{Mitte} des Zeitfensters \\
\bottomrule
\end{tabular}
\end{center}

Die Standardspannung betraegt \texttt{+1 V} fuer Impulse und \texttt{0 V}
fuer den Ruhepegel. Zeitverhalten und Amplitude sind vollstaendig in
\texttt{config.py} konfigurierbar.

% ------------------------------------------------------------------
\subsection*{Hardwareanforderungen}

\begin{itemize}
  \item \textbf{Digilent Analog Discovery 2} (AD2) -- USB-Multifunktionsinstrument
  \begin{itemize}
    \item Wird als Arbitrary Waveform Generator (AWG) verwendet, um das
          PPM-Signal am \textbf{Analogausgang Kanal 1 (W1)} auszugeben
  \end{itemize}
  \item Ein USB-Kabel zur Verbindung des AD2 mit dem PC
\end{itemize}

% ------------------------------------------------------------------
\subsection*{Softwarevoraussetzungen}

\subsubsection*{1. WaveForms (Digilent SDK)}

Das Programm benoetigt die \textbf{DWF-Bibliothek (Device WaveForms)}
(\texttt{dwf.dll} unter Windows, \texttt{libdwf.so} unter Linux,
\texttt{libdwf.dylib} unter macOS).

WaveForms kann von der Digilent-Website heruntergeladen und installiert werden:\\
\url{https://digilent.com/reference/software/waveforms/waveforms-3/start}

\textbf{Wichtig:} Die WaveForms-Desktopanwendung muss geschlossen sein, bevor
ein Python-Skript ausgefuehrt wird, da WaveForms und ein Python-Skript das
Geraet nicht gleichzeitig steuern koennen.

\subsubsection*{2. Python 3}

Python 3 installieren (Version 3.7 oder hoeher empfohlen).

\subsubsection*{3. Python-Pakete}

Das benoetigt Paket installieren:

\begin{lstlisting}[language=bash]
pip install matplotlib
\end{lstlisting}

\texttt{ctypes} und \texttt{time} sind Bestandteil der Python-Standardbibliothek
und muessen nicht separat installiert werden.

% ------------------------------------------------------------------
\subsection*{Dateistruktur}

\begin{lstlisting}
VLC_Waveform SDK/
+-- main_PPM.py       <- Hauptskript -- diese Datei ausfuehren
+-- PPM.py            <- Logik zur PPM-Signalerzeugung
+-- config.py         <- Alle konfigurierbaren Parameter und Geraeteeinrichtung
+-- wavegen.py        <- Low-Level-Hilfsfunktionen fuer den Waveformgenerator
+-- device.py         <- Hilfsfunktionen fuer Geraeteverbindung und Fehlerbehandlung
+-- dwfconstants.py   <- DWF-Bibliothekskonstanten
+-- main.py           <- Eigenstaendige Demo: Rechteckwelle (separater Test)
+-- waveform_data.csv <- (optional) Exportierte Wellenformdaten
\end{lstlisting}

% ------------------------------------------------------------------
\subsection*{Konfiguration}

Vor dem Starten des Programms sollten die Parameter in \textbf{\texttt{config.py}}
geprueft und bei Bedarf angepasst werden:

\begin{center}
\begin{tabular}{lll}
\toprule
\textbf{Parameter} & \textbf{Standardwert} & \textbf{Beschreibung} \\
\midrule
\texttt{DESIRED\_PEAK\_VOLTAGE} & \texttt{1.0} V & Impulsamplitude \\
\texttt{REST\_LEVEL} & \texttt{0.0} V & Spannung waehrend des Ruhepegels \\
\texttt{BIT\_DURATION} & \texttt{1e-6} s (1\,\textmu s) & Dauer eines Bit-Zeitfensters \\
\texttt{PULSE\_WIDTH} & \texttt{BIT\_DURATION / 2} & Impulsbreite innerhalb eines Zeitfensters \\
\texttt{SAMPLES\_PER\_BIT\_GEN} & \texttt{2048} & AWG-Abtastwerte pro Bit \\
\texttt{LOOP\_CONTINUOUSLY} & \texttt{True} & Wellenform wiederholen (\texttt{True}) oder einmalig (\texttt{False}) \\
\texttt{CHANNEL\_AWG} & \texttt{0} & AWG-Ausgangskanal (0 = W1) \\
\texttt{save\_to\_CSV} & \texttt{False} & Abtastwerte in \texttt{waveform\_data.csv} speichern \\
\texttt{diagnostic} & \texttt{False} & Signaldetails in der Konsole ausgeben \\
\bottomrule
\end{tabular}
\end{center}

% ------------------------------------------------------------------
\subsection*{Programm starten}

\begin{enumerate}
  \item Den Analog Discovery 2 per USB mit dem Computer verbinden.
  \item Die WaveForms-Desktopanwendung schliessen, falls sie geoeffnet ist.
  \item Ein Terminal oeffnen und in den SDK-Ordner wechseln:
\begin{lstlisting}[language=bash]
cd "VLC_Waveform SDK"
\end{lstlisting}
  \item Das Hauptskript ausfuehren:
\begin{lstlisting}[language=bash]
python main_PPM.py
\end{lstlisting}
  Bei Erfolg erscheint folgende Ausgabe:
\begin{lstlisting}
Device opened successfully.
Bitfolge in Zahlen eintragen...
\end{lstlisting}
\end{enumerate}

% ------------------------------------------------------------------
\subsection*{Programm verwenden}

Nach dem Start wechselt das Programm in eine interaktive Schleife:

\begin{enumerate}
  \item Eine binaere Bitfolge eingeben, wenn der Prompt erscheint --
        eine Folge aus \texttt{0} und \texttt{1}, zum Beispiel:
\begin{lstlisting}
Bitfolge in Zahlen eintragen...01101001
\end{lstlisting}
  \item Das Programm fuehrt anschliessend folgende Schritte aus:
  \begin{itemize}
    \item Es erzeugt die entsprechende PPM-Wellenform
    \item Es gibt das Signal ueber den AWG-Kanal (W1) des Analog Discovery 2 aus
    \item Es wartet, bis die Signaldauer abgelaufen ist, und fragt dann erneut
          nach einer Eingabe
  \end{itemize}
  \item Den Vorgang mit einer neuen Bitfolge beliebig oft wiederholen.
  \item Das Programm mit \texttt{Ctrl+C} beenden. Der AWG-Ausgang wird gestoppt
        und das Geraet ordnungsgemaess getrennt.
\end{enumerate}

\textbf{Hinweis:} Das erzeugte Signal kann mit dem Oszilloskopeingang des
Analog Discovery 2 (Kanal 1 / 1+) und dem WaveForms-Scope-Instrument beobachtet
werden -- jedoch erst nach dem Beenden von \texttt{main\_PPM.py}, da jeweils nur
eine Anwendung das Geraet steuern kann.

% ------------------------------------------------------------------
\subsection*{Fehlerbehebung}

\begin{center}
\begin{tabular}{p{0.38\textwidth}p{0.55\textwidth}}
\toprule
\textbf{Problem} & \textbf{Loesung} \\
\midrule
\texttt{Failed to open device} &
  Sicherstellen, dass der AD2 per USB verbunden und WaveForms geschlossen ist \\
\addlinespace
\texttt{Could not load dwf library} &
  WaveForms von Digilent installieren; sicherstellen, dass die DWF-Bibliothek
  im System-PATH vorhanden ist \\
\addlinespace
\texttt{No device detected} &
  Einen anderen USB-Anschluss oder ein anderes Kabel verwenden;
  Geraetemanager bzw.\ \texttt{lsusb} pruefen \\
\addlinespace
Skript haengt oder keine Ausgabe &
  \texttt{BIT\_DURATION} und Bitfolgenlange pruefen; diese sollten nicht
  zu gross sein \\
\addlinespace
\texttt{ModuleNotFoundError: matplotlib} &
  \texttt{pip install matplotlib} ausfuehren \\
\bottomrule
\end{tabular}
\end{center}
 in ein LaTeX-Hauptdokument eingebunden werden.
%
% Benoetigt folgende Pakete im Hauptdokument:
%   \usepackage[T1]{fontenc}
%   \usepackage[utf8]{inputenc}
%   \usepackage{booktabs}
%   \usepackage{listings}
%   \usepackage{hyperref}
%   \usepackage{enumitem}
%   \usepackage{microtype}

\section*{VLC -- PPM-Signalgenerator (Analog Discovery 2)}

Ein Python-basierter Pulspositionsmodulations-Signalgenerator (PPM),
der den \textbf{Digilent Analog Discovery 2} als Hardware-Ausgabegeraet verwendet.

% ------------------------------------------------------------------
\subsection*{Uebersicht}

Dieses Programm erzeugt ein PPM-Signal (Pulspositionsmodulation) aus einer
vom Benutzer eingegebenen binaeren Bitfolge und gibt es in Echtzeit ueber den
\textbf{Arbitrary Waveform Generator (AWG)} des Analog Discovery 2 aus.

Der Einstiegspunkt des Programms ist \textbf{\texttt{main\_PPM.py}} im
Verzeichnis \texttt{VLC\_Waveform SDK/}.

% ------------------------------------------------------------------
\subsection*{Funktionsweise von PPM}

Bei PPM wird jedes Bit durch einen Impuls dargestellt, der an einer bestimmten
Position innerhalb eines festen Zeitfensters platziert wird:

\begin{center}
\begin{tabular}{ll}
\toprule
\textbf{Bit} & \textbf{Impulsposition} \\
\midrule
\texttt{0} & Impuls am \textbf{Anfang} des Zeitfensters \\
\texttt{1} & Impuls in der \textbf{Mitte} des Zeitfensters \\
\bottomrule
\end{tabular}
\end{center}

Die Standardspannung betraegt \texttt{+1 V} fuer Impulse und \texttt{0 V}
fuer den Ruhepegel. Zeitverhalten und Amplitude sind vollstaendig in
\texttt{config.py} konfigurierbar.

% ------------------------------------------------------------------
\subsection*{Hardwareanforderungen}

\begin{itemize}
  \item \textbf{Digilent Analog Discovery 2} (AD2) -- USB-Multifunktionsinstrument
  \begin{itemize}
    \item Wird als Arbitrary Waveform Generator (AWG) verwendet, um das
          PPM-Signal am \textbf{Analogausgang Kanal 1 (W1)} auszugeben
  \end{itemize}
  \item Ein USB-Kabel zur Verbindung des AD2 mit dem PC
\end{itemize}

% ------------------------------------------------------------------
\subsection*{Softwarevoraussetzungen}

\subsubsection*{1. WaveForms (Digilent SDK)}

Das Programm benoetigt die \textbf{DWF-Bibliothek (Device WaveForms)}
(\texttt{dwf.dll} unter Windows, \texttt{libdwf.so} unter Linux,
\texttt{libdwf.dylib} unter macOS).

WaveForms kann von der Digilent-Website heruntergeladen und installiert werden:\\
\url{https://digilent.com/reference/software/waveforms/waveforms-3/start}

\textbf{Wichtig:} Die WaveForms-Desktopanwendung muss geschlossen sein, bevor
ein Python-Skript ausgefuehrt wird, da WaveForms und ein Python-Skript das
Geraet nicht gleichzeitig steuern koennen.

\subsubsection*{2. Python 3}

Python 3 installieren (Version 3.7 oder hoeher empfohlen).

\subsubsection*{3. Python-Pakete}

Das benoetigt Paket installieren:

\begin{lstlisting}[language=bash]
pip install matplotlib
\end{lstlisting}

\texttt{ctypes} und \texttt{time} sind Bestandteil der Python-Standardbibliothek
und muessen nicht separat installiert werden.

% ------------------------------------------------------------------
\subsection*{Dateistruktur}

\begin{lstlisting}
VLC_Waveform SDK/
+-- main_PPM.py       <- Hauptskript -- diese Datei ausfuehren
+-- PPM.py            <- Logik zur PPM-Signalerzeugung
+-- config.py         <- Alle konfigurierbaren Parameter und Geraeteeinrichtung
+-- wavegen.py        <- Low-Level-Hilfsfunktionen fuer den Waveformgenerator
+-- device.py         <- Hilfsfunktionen fuer Geraeteverbindung und Fehlerbehandlung
+-- dwfconstants.py   <- DWF-Bibliothekskonstanten
+-- main.py           <- Eigenstaendige Demo: Rechteckwelle (separater Test)
+-- waveform_data.csv <- (optional) Exportierte Wellenformdaten
\end{lstlisting}

% ------------------------------------------------------------------
\subsection*{Konfiguration}

Vor dem Starten des Programms sollten die Parameter in \textbf{\texttt{config.py}}
geprueft und bei Bedarf angepasst werden:

\begin{center}
\begin{tabular}{lll}
\toprule
\textbf{Parameter} & \textbf{Standardwert} & \textbf{Beschreibung} \\
\midrule
\texttt{DESIRED\_PEAK\_VOLTAGE} & \texttt{1.0} V & Impulsamplitude \\
\texttt{REST\_LEVEL} & \texttt{0.0} V & Spannung waehrend des Ruhepegels \\
\texttt{BIT\_DURATION} & \texttt{1e-6} s (1\,\textmu s) & Dauer eines Bit-Zeitfensters \\
\texttt{PULSE\_WIDTH} & \texttt{BIT\_DURATION / 2} & Impulsbreite innerhalb eines Zeitfensters \\
\texttt{SAMPLES\_PER\_BIT\_GEN} & \texttt{2048} & AWG-Abtastwerte pro Bit \\
\texttt{LOOP\_CONTINUOUSLY} & \texttt{True} & Wellenform wiederholen (\texttt{True}) oder einmalig (\texttt{False}) \\
\texttt{CHANNEL\_AWG} & \texttt{0} & AWG-Ausgangskanal (0 = W1) \\
\texttt{save\_to\_CSV} & \texttt{False} & Abtastwerte in \texttt{waveform\_data.csv} speichern \\
\texttt{diagnostic} & \texttt{False} & Signaldetails in der Konsole ausgeben \\
\bottomrule
\end{tabular}
\end{center}

% ------------------------------------------------------------------
\subsection*{Programm starten}

\begin{enumerate}
  \item Den Analog Discovery 2 per USB mit dem Computer verbinden.
  \item Die WaveForms-Desktopanwendung schliessen, falls sie geoeffnet ist.
  \item Ein Terminal oeffnen und in den SDK-Ordner wechseln:
\begin{lstlisting}[language=bash]
cd "VLC_Waveform SDK"
\end{lstlisting}
  \item Das Hauptskript ausfuehren:
\begin{lstlisting}[language=bash]
python main_PPM.py
\end{lstlisting}
  Bei Erfolg erscheint folgende Ausgabe:
\begin{lstlisting}
Device opened successfully.
Bitfolge in Zahlen eintragen...
\end{lstlisting}
\end{enumerate}

% ------------------------------------------------------------------
\subsection*{Programm verwenden}

Nach dem Start wechselt das Programm in eine interaktive Schleife:

\begin{enumerate}
  \item Eine binaere Bitfolge eingeben, wenn der Prompt erscheint --
        eine Folge aus \texttt{0} und \texttt{1}, zum Beispiel:
\begin{lstlisting}
Bitfolge in Zahlen eintragen...01101001
\end{lstlisting}
  \item Das Programm fuehrt anschliessend folgende Schritte aus:
  \begin{itemize}
    \item Es erzeugt die entsprechende PPM-Wellenform
    \item Es gibt das Signal ueber den AWG-Kanal (W1) des Analog Discovery 2 aus
    \item Es wartet, bis die Signaldauer abgelaufen ist, und fragt dann erneut
          nach einer Eingabe
  \end{itemize}
  \item Den Vorgang mit einer neuen Bitfolge beliebig oft wiederholen.
  \item Das Programm mit \texttt{Ctrl+C} beenden. Der AWG-Ausgang wird gestoppt
        und das Geraet ordnungsgemaess getrennt.
\end{enumerate}

\textbf{Hinweis:} Das erzeugte Signal kann mit dem Oszilloskopeingang des
Analog Discovery 2 (Kanal 1 / 1+) und dem WaveForms-Scope-Instrument beobachtet
werden -- jedoch erst nach dem Beenden von \texttt{main\_PPM.py}, da jeweils nur
eine Anwendung das Geraet steuern kann.

% ------------------------------------------------------------------
\subsection*{Fehlerbehebung}

\begin{center}
\begin{tabular}{p{0.38\textwidth}p{0.55\textwidth}}
\toprule
\textbf{Problem} & \textbf{Loesung} \\
\midrule
\texttt{Failed to open device} &
  Sicherstellen, dass der AD2 per USB verbunden und WaveForms geschlossen ist \\
\addlinespace
\texttt{Could not load dwf library} &
  WaveForms von Digilent installieren; sicherstellen, dass die DWF-Bibliothek
  im System-PATH vorhanden ist \\
\addlinespace
\texttt{No device detected} &
  Einen anderen USB-Anschluss oder ein anderes Kabel verwenden;
  Geraetemanager bzw.\ \texttt{lsusb} pruefen \\
\addlinespace
Skript haengt oder keine Ausgabe &
  \texttt{BIT\_DURATION} und Bitfolgenlange pruefen; diese sollten nicht
  zu gross sein \\
\addlinespace
\texttt{ModuleNotFoundError: matplotlib} &
  \texttt{pip install matplotlib} ausfuehren \\
\bottomrule
\end{tabular}
\end{center}
 in ein LaTeX-Hauptdokument eingebunden werden.
%
% Benoetigt folgende Pakete im Hauptdokument:
%   \usepackage[T1]{fontenc}
%   \usepackage[utf8]{inputenc}
%   \usepackage{booktabs}
%   \usepackage{listings}
%   \usepackage{hyperref}
%   \usepackage{enumitem}
%   \usepackage{microtype}

\section*{VLC -- PPM-Signalgenerator (Analog Discovery 2)}

Ein Python-basierter Pulspositionsmodulations-Signalgenerator (PPM),
der den \textbf{Digilent Analog Discovery 2} als Hardware-Ausgabegeraet verwendet.

% ------------------------------------------------------------------
\subsection*{Uebersicht}

Dieses Programm erzeugt ein PPM-Signal (Pulspositionsmodulation) aus einer
vom Benutzer eingegebenen binaeren Bitfolge und gibt es in Echtzeit ueber den
\textbf{Arbitrary Waveform Generator (AWG)} des Analog Discovery 2 aus.

Der Einstiegspunkt des Programms ist \textbf{\texttt{main\_PPM.py}} im
Verzeichnis \texttt{VLC\_Waveform SDK/}.

% ------------------------------------------------------------------
\subsection*{Funktionsweise von PPM}

Bei PPM wird jedes Bit durch einen Impuls dargestellt, der an einer bestimmten
Position innerhalb eines festen Zeitfensters platziert wird:

\begin{center}
\begin{tabular}{ll}
\toprule
\textbf{Bit} & \textbf{Impulsposition} \\
\midrule
\texttt{0} & Impuls am \textbf{Anfang} des Zeitfensters \\
\texttt{1} & Impuls in der \textbf{Mitte} des Zeitfensters \\
\bottomrule
\end{tabular}
\end{center}

Die Standardspannung betraegt \texttt{+1 V} fuer Impulse und \texttt{0 V}
fuer den Ruhepegel. Zeitverhalten und Amplitude sind vollstaendig in
\texttt{config.py} konfigurierbar.

% ------------------------------------------------------------------
\subsection*{Hardwareanforderungen}

\begin{itemize}
  \item \textbf{Digilent Analog Discovery 2} (AD2) -- USB-Multifunktionsinstrument
  \begin{itemize}
    \item Wird als Arbitrary Waveform Generator (AWG) verwendet, um das
          PPM-Signal am \textbf{Analogausgang Kanal 1 (W1)} auszugeben
  \end{itemize}
  \item Ein USB-Kabel zur Verbindung des AD2 mit dem PC
\end{itemize}

% ------------------------------------------------------------------
\subsection*{Softwarevoraussetzungen}

\subsubsection*{1. WaveForms (Digilent SDK)}

Das Programm benoetigt die \textbf{DWF-Bibliothek (Device WaveForms)}
(\texttt{dwf.dll} unter Windows, \texttt{libdwf.so} unter Linux,
\texttt{libdwf.dylib} unter macOS).

WaveForms kann von der Digilent-Website heruntergeladen und installiert werden:\\
\url{https://digilent.com/reference/software/waveforms/waveforms-3/start}

\textbf{Wichtig:} Die WaveForms-Desktopanwendung muss geschlossen sein, bevor
ein Python-Skript ausgefuehrt wird, da WaveForms und ein Python-Skript das
Geraet nicht gleichzeitig steuern koennen.

\subsubsection*{2. Python 3}

Python 3 installieren (Version 3.7 oder hoeher empfohlen).

\subsubsection*{3. Python-Pakete}

Das benoetigt Paket installieren:

\begin{lstlisting}[language=bash]
pip install matplotlib
\end{lstlisting}

\texttt{ctypes} und \texttt{time} sind Bestandteil der Python-Standardbibliothek
und muessen nicht separat installiert werden.

% ------------------------------------------------------------------
\subsection*{Dateistruktur}

\begin{lstlisting}
VLC_Waveform SDK/
+-- main_PPM.py       <- Hauptskript -- diese Datei ausfuehren
+-- PPM.py            <- Logik zur PPM-Signalerzeugung
+-- config.py         <- Alle konfigurierbaren Parameter und Geraeteeinrichtung
+-- wavegen.py        <- Low-Level-Hilfsfunktionen fuer den Waveformgenerator
+-- device.py         <- Hilfsfunktionen fuer Geraeteverbindung und Fehlerbehandlung
+-- dwfconstants.py   <- DWF-Bibliothekskonstanten
+-- main.py           <- Eigenstaendige Demo: Rechteckwelle (separater Test)
+-- waveform_data.csv <- (optional) Exportierte Wellenformdaten
\end{lstlisting}

% ------------------------------------------------------------------
\subsection*{Konfiguration}

Vor dem Starten des Programms sollten die Parameter in \textbf{\texttt{config.py}}
geprueft und bei Bedarf angepasst werden:

\begin{center}
\begin{tabular}{lll}
\toprule
\textbf{Parameter} & \textbf{Standardwert} & \textbf{Beschreibung} \\
\midrule
\texttt{DESIRED\_PEAK\_VOLTAGE} & \texttt{1.0} V & Impulsamplitude \\
\texttt{REST\_LEVEL} & \texttt{0.0} V & Spannung waehrend des Ruhepegels \\
\texttt{BIT\_DURATION} & \texttt{1e-6} s (1\,\textmu s) & Dauer eines Bit-Zeitfensters \\
\texttt{PULSE\_WIDTH} & \texttt{BIT\_DURATION / 2} & Impulsbreite innerhalb eines Zeitfensters \\
\texttt{SAMPLES\_PER\_BIT\_GEN} & \texttt{2048} & AWG-Abtastwerte pro Bit \\
\texttt{LOOP\_CONTINUOUSLY} & \texttt{True} & Wellenform wiederholen (\texttt{True}) oder einmalig (\texttt{False}) \\
\texttt{CHANNEL\_AWG} & \texttt{0} & AWG-Ausgangskanal (0 = W1) \\
\texttt{save\_to\_CSV} & \texttt{False} & Abtastwerte in \texttt{waveform\_data.csv} speichern \\
\texttt{diagnostic} & \texttt{False} & Signaldetails in der Konsole ausgeben \\
\bottomrule
\end{tabular}
\end{center}

% ------------------------------------------------------------------
\subsection*{Programm starten}

\begin{enumerate}
  \item Den Analog Discovery 2 per USB mit dem Computer verbinden.
  \item Die WaveForms-Desktopanwendung schliessen, falls sie geoeffnet ist.
  \item Ein Terminal oeffnen und in den SDK-Ordner wechseln:
\begin{lstlisting}[language=bash]
cd "VLC_Waveform SDK"
\end{lstlisting}
  \item Das Hauptskript ausfuehren:
\begin{lstlisting}[language=bash]
python main_PPM.py
\end{lstlisting}
  Bei Erfolg erscheint folgende Ausgabe:
\begin{lstlisting}
Device opened successfully.
Bitfolge in Zahlen eintragen...
\end{lstlisting}
\end{enumerate}

% ------------------------------------------------------------------
\subsection*{Programm verwenden}

Nach dem Start wechselt das Programm in eine interaktive Schleife:

\begin{enumerate}
  \item Eine binaere Bitfolge eingeben, wenn der Prompt erscheint --
        eine Folge aus \texttt{0} und \texttt{1}, zum Beispiel:
\begin{lstlisting}
Bitfolge in Zahlen eintragen...01101001
\end{lstlisting}
  \item Das Programm fuehrt anschliessend folgende Schritte aus:
  \begin{itemize}
    \item Es erzeugt die entsprechende PPM-Wellenform
    \item Es gibt das Signal ueber den AWG-Kanal (W1) des Analog Discovery 2 aus
    \item Es wartet, bis die Signaldauer abgelaufen ist, und fragt dann erneut
          nach einer Eingabe
  \end{itemize}
  \item Den Vorgang mit einer neuen Bitfolge beliebig oft wiederholen.
  \item Das Programm mit \texttt{Ctrl+C} beenden. Der AWG-Ausgang wird gestoppt
        und das Geraet ordnungsgemaess getrennt.
\end{enumerate}

\textbf{Hinweis:} Das erzeugte Signal kann mit dem Oszilloskopeingang des
Analog Discovery 2 (Kanal 1 / 1+) und dem WaveForms-Scope-Instrument beobachtet
werden -- jedoch erst nach dem Beenden von \texttt{main\_PPM.py}, da jeweils nur
eine Anwendung das Geraet steuern kann.

% ------------------------------------------------------------------
\subsection*{Fehlerbehebung}

\begin{center}
\begin{tabular}{p{0.38\textwidth}p{0.55\textwidth}}
\toprule
\textbf{Problem} & \textbf{Loesung} \\
\midrule
\texttt{Failed to open device} &
  Sicherstellen, dass der AD2 per USB verbunden und WaveForms geschlossen ist \\
\addlinespace
\texttt{Could not load dwf library} &
  WaveForms von Digilent installieren; sicherstellen, dass die DWF-Bibliothek
  im System-PATH vorhanden ist \\
\addlinespace
\texttt{No device detected} &
  Einen anderen USB-Anschluss oder ein anderes Kabel verwenden;
  Geraetemanager bzw.\ \texttt{lsusb} pruefen \\
\addlinespace
Skript haengt oder keine Ausgabe &
  \texttt{BIT\_DURATION} und Bitfolgenlange pruefen; diese sollten nicht
  zu gross sein \\
\addlinespace
\texttt{ModuleNotFoundError: matplotlib} &
  \texttt{pip install matplotlib} ausfuehren \\
\bottomrule
\end{tabular}
\end{center}
